% ============================================================================
% Fichier  : partie2/P02_C11_conclure_rapport.tex
% Titre    : Conclure et Rédiger : Mesures, Quadrant, Observation, Rapport
% Auteur   : MINKA MI NGUIDJOÏ Thierry Emmanuel
% Version  : 1.0 -- novembre 2026
% ============================================================================
% Base : ACSI C09 (mesures, quadrant, seuils µ 0,5 et U 0,20, règle de
% décision, cinq domaines de LiZbiZ A 0,80/0,05/0,88 ; B 0,67/0,08/1,19 ;
% C 0,08/0,08/0,82 ; D 0,65/0,25/1,24 ; E 0,40/0,30/1,57 ; observation
% D.2.1 en sept éléments) et ACSI C22 (décomposition de l'entropie,
% D 66 %, E 56 %, B 35 %).
% Sources : A1 (hal-05701270) pour les mesures et le quadrant ; A0 pour
% la non-conformité connue et inconnue ; Shannon (1948) et Cover et Thomas
% (2006) pour la règle de chaînage de l'entropie. L'identité
% H = H_b(U) + (1-U) H_b(µ/(1-U)) est donnée exacte ; la Proposition 1
% d'A1 comporte un terme en trop (erratum signalé à l'auteur), elle n'est
% donc pas citée pour cette identité.
% Corrigés : corriges/Corr_C11.tex.
% ============================================================================

\chapter[Conclure et Rédiger]{Conclure et Rédiger : Mesures, Quadrant, Observation, Rapport}
\label{chap:conclure}
\chaptermark{Conclure et rédiger}

\epigraph{%
    « Conclure, ce n'est pas compter les \emph{oui}.\\
    C'est savoir ce qu'ils valent. »%
}{Principe de la méthode}

\niveauChapitre{Intermédiaire}{%
    Maîtrise du Chapitre~\ref{chap:preuve} (instrument \OUI{}, \NON{},
    \NEK{}, preuves, papiers de travail)%
}

\begin{boxobjectifs}
\begin{itemize}
    \item Calculer, pour une question et pour un domaine, le taux de
          conformité~$\mu$, le taux d'incertitude~$U$ et l'entropie~$H$.
    \item Démontrer la décomposition de l'entropie en une part
          d'ignorance et une part de désaccord, et l'interpréter.
    \item Situer un domaine dans le quadrant conformité × incertitude et en
          déduire l'action appropriée.
    \item Appliquer la règle de décision du cours, qui combine toujours
          $\mu$, $U$ et $H$.
    \item Graduer et rédiger une observation en sept éléments.
    \item Construire un rapport d'audit complet, de la synthèse aux
          annexes.
\end{itemize}
\end{boxobjectifs}

\bigskip

% ============================================================================
\section*{Question directrice : que peut-on affirmer~?}
% ============================================================================

Le terrain est terminé. AUDIX CONSEIL dispose, pour les cinq domaines de sa
mission chez LiZbiZ, de centaines de réponses et de dizaines de tests. Le
directeur général attend une phrase~: « Mes contrôles fonctionnent-ils~? »
M.~BEYALA sait qu'il doit répondre domaine par domaine, en disant à la fois
ce qui va mal et ce que l'on ne sait pas.

\medskip
\textbf{Plan du chapitre.}
\begin{enumerate}
    \item Les trois mesures (section~\ref{sec:c11-mesures}).
    \item La décomposition de l'entropie (section~\ref{sec:c11-decomp}).
    \item Le quadrant (section~\ref{sec:c11-quadrant}).
    \item La règle de décision (section~\ref{sec:c11-regle}).
    \item L'observation en sept éléments (section~\ref{sec:c11-observation}).
    \item Le rapport (section~\ref{sec:c11-rapport}).
\end{enumerate}
Le suivi, la comparaison des cycles et la maturité font l'objet du
Chapitre~\ref{chap:suivi}.

\separateur

% ============================================================================
\section{Trois Mesures}
\label{sec:c11-mesures}
% ============================================================================

\subsection{Pour une question}

Pour une question~$q$ à laquelle plusieurs répondants ont répondu, on note
$p_q(\OUI)$, $p_q(\NON)$ et $p_q(\NEK)$ les parts de chaque réponse. L'article
A1 définit trois mesures~\cite{minka2026nek}.

\begin{boxdef}
\termecle{Taux de conformité}\index{taux de conformité}~: $\mu_q =
p_q(\OUI)$, la part des réponses \OUI{}.

\medskip
\termecle{Taux d'incertitude}\index{taux d'incertitude}~: $U_q = p_q(\NEK)$,
la part des réponses \NEK{}.

\medskip
\termecle{Entropie}\index{entropie}~: dispersion des réponses, en bits,
\[
    H_q = -\sum_{m \in \{\OUI, \NON, \NEK\}} p_q(m) \log_2 p_q(m),
    \qquad 0 \leq H_q \leq \log_2 3 \approx 1{,}585,
\]
avec la convention $0 \log_2 0 = 0$.
\end{boxdef}

L'entropie, introduite par Claude Shannon en 1948~\cite{shannon1948},
mesure la quantité d'information qu'apporterait la connaissance de la
réponse~: elle est nulle quand tout le monde répond la même chose, maximale
quand les trois réponses sont également fréquentes.

\subsection{Pour un domaine}

Au niveau d'un domaine ou d'un sous-domaine, le cours agrège les réponses de
toutes ses questions avant de calculer les trois mesures. C'est ce qui a été
fait pour le domaine~D au prologue (20 réponses) et au
Chapitre~\ref{chap:decomposition} (sous-domaine par sous-domaine). Les
mesures d'un niveau se recalculent toujours à partir des réponses, jamais
comme moyenne des mesures du niveau inférieur.

\subsection{Quatre cas limites}

\begin{tableaucours}{L{4.6cm}cccL{3.4cm}}{%
    \tableauHeader{Situation} &
    \tableauHeader{$\mu$} & \tableauHeader{$U$} & \tableauHeader{$H$} &
    \tableauHeader{Lecture}
}
Tous répondent \OUI{} & 1 & 0 & 0 & Conformité certaine \\
Tous répondent \NON{} & 0 & 0 & 0 & Non-conformité certaine \\
Tous répondent \NEK{} & 0 & 1 & 0 & \textbf{Ignorance unanime} \\
Réponses également réparties & $1/3$ & $1/3$ & 1,585 & Dispersion maximale \\
\end{tableaucours}
\vspace{-0.8em}
\begin{center}{\small\itshape Tableau~11.1. Les cas limites des trois
mesures (d'après~\cite{minka2026nek})}\end{center}

\begin{boxalert}
\textbf{Violation typique.} Conclure «~certitude~» parce que l'entropie est
nulle.

\smallskip
\textbf{Règle~:} une entropie nulle signifie seulement que tous les
répondants ont donné la même réponse. Si cette réponse est \NEK{}, personne
ne sait rien~: $H = 0$ et $U = 1$. On ne lit jamais $H$ sans $U$ ni $\mu$.
\end{boxalert}

\begin{boxoutils}{Calculer les trois mesures}
\begin{lstlisting}[language=Python]
from math import log2
def mesures(oui, non, nek):
    n = oui + non + nek
    p = [oui / n, non / n, nek / n]
    H = -sum(x * log2(x) for x in p if x > 0)
    return round(oui / n, 2), round(nek / n, 2), round(H, 2)

print(mesures(13, 2, 5))   # domaine D : (0.65, 0.25, 1.24)
\end{lstlisting}
\end{boxoutils}

% ============================================================================
\section{Décomposer l'Entropie}
\label{sec:c11-decomp}
% ============================================================================

\subsection{Deux sources de dispersion}

Une entropie élevée peut venir de deux choses très différentes~: beaucoup
de répondants ne savent pas (ignorance), ou ceux qui savent ne sont pas
d'accord (désaccord entre \OUI{} et \NON{}). L'auditeur doit savoir
laquelle domine~: on ne traite pas l'ignorance comme le désaccord.

\subsection{La démonstration}

On regroupe d'abord les réponses en deux classes~: \NEK{} (probabilité $U$)
et « réponse informée » (\OUI{} ou \NON{}, probabilité $1 - U$). À
l'intérieur de la seconde classe, la part des \OUI{} vaut $\mu / (1 - U)$.
La règle de chaînage de l'entropie~\cite{shannon1948,cover2006elements}
donne alors, pour $U < 1$~:
\[
    H = H_b(U) + (1 - U)\, H_b\!\left(\frac{\mu}{1 - U}\right),
\]
où $H_b(x) = -x\log_2 x - (1-x)\log_2(1-x)$ est l'entropie binaire.

\begin{boxinfo}
\textbf{Vérification directe.} Avec $p_{\NEK} = U$, $p_{\OUI} = \mu$ et
$p_{\NON} = 1 - U - \mu$, on développe le second membre~:
\begin{align*}
H_b(U) &= -U\log_2 U - (1-U)\log_2(1-U),\\
(1-U)\,H_b\!\left(\tfrac{\mu}{1-U}\right)
 &= -\mu\log_2\tfrac{\mu}{1-U} - (1-U-\mu)\log_2\tfrac{1-U-\mu}{1-U}.
\end{align*}
En additionnant, les termes en $\log_2(1-U)$ se compensent
($-(1-U) + \mu + (1-U-\mu) = 0$), et il reste
$-U\log_2 U - \mu\log_2\mu - (1-U-\mu)\log_2(1-U-\mu)$, c'est-à-dire $H$.
\end{boxinfo}

Le premier terme, $H_b(U)$, est la \textbf{part d'ignorance}~: il ne
dépend que de la proportion de \NEK{}. Le second est la \textbf{part de
désaccord}~: la dispersion entre \OUI{} et \NON{} chez ceux qui savent,
pondérée par leur proportion.

\subsection{Les domaines de LiZbiZ}

\begin{tableaucours}{L{3.0cm}ccccc}{%
    \tableauHeader{Domaine} &
    \tableauHeader{$H$} &
    \tableauHeader{$H_b(U)$} &
    \tableauHeader{Désaccord} &
    \tableauHeader{Part d'ignorance} &
    \tableauHeader{Lecture}
}
B. Sauvegardes & 1,19 & 0,41 & 0,78 & 35~\% & Surtout désaccord \\
D. Contrôle d'accès & 1,24 & 0,81 & 0,43 & 66~\% & Surtout ignorance \\
E. Exploitation & 1,57 & 0,88 & 0,69 & 56~\% & Les deux \\
\end{tableaucours}
\vspace{-0.8em}
\begin{center}{\small\itshape Tableau~11.2. Décomposition de l'entropie de
trois domaines}\end{center}

\begin{boxlecture}{Deux entropies voisines}
B et D ont des entropies voisines (1,19 et 1,24), mais pour des raisons
opposées.

\medskip
\textbf{Analyse} (confiance élevée). Dans le domaine~B, deux tiers de
l'entropie viennent du désaccord entre \OUI{} et \NON{}~: les répondants
savent, mais les sites diffèrent (Garoua ne teste pas ses restaurations).
On décompose par site. Dans le domaine~D, deux tiers viennent de
l'ignorance~: on cherche la preuve. Le même chiffre appelle deux
démarches.

\medskip
\hfill\lecture{0{,}65}{0{,}25}{1{,}24}
\end{boxlecture}

% ============================================================================
\section{Le Quadrant Conformité × Incertitude}
\label{sec:c11-quadrant}
% ============================================================================

L'article A1 propose de lire ensemble la conformité et l'incertitude dans un
quadrant à quatre régions, chacune appelant une action
différente~\cite{minka2026nek}. Le cours retient, par convention, un seuil de
0,5 pour $\mu$ et de 0,20 pour $U$~; chaque mission peut fixer les siens, à
condition de les écrire dans le plan d'audit.

\noindent\begin{minipage}{\linewidth}
\begin{center}
\begin{tikzpicture}[scale=1.0]
\fill[green!18] (3,0) rectangle (6,2);
\fill[red!15] (0,0) rectangle (3,2);
\fill[orange!25] (3,2) rectangle (6,5);
\fill[orange!15] (0,2) rectangle (3,5);
\draw[thick, color=CouleurPrimaire] (0,0) rectangle (6,5);
\draw[dashed, color=CouleurPrimaire] (3,0) -- (3,5);
\draw[dashed, color=CouleurPrimaire] (0,2) -- (6,2);
\node[align=center, font=\footnotesize\bfseries] at (1.5,1.0) {Non-conformité\\ claire};
\node[align=center, font=\footnotesize\bfseries] at (4.5,1.0) {Conformité\\ claire};
\node[align=center, font=\footnotesize\bfseries] at (1.5,3.5) {Non-conformité\\ opaque};
\node[align=center, font=\footnotesize\bfseries] at (4.5,3.5) {Conformité\\ opaque};
\foreach \x/\y/\t in {4.8/0.5/A, 4.0/0.8/B, 0.5/0.8/C, 3.9/2.5/D, 2.4/3.0/E}
    \fill[CouleurAccent] (\x,\y) circle (2.5pt) node[above right, font=\scriptsize\bfseries] {\t};
\node[font=\scriptsize] at (3,-0.3) {0,5};
\node[font=\scriptsize] at (-0.35,2) {0,20};
\node[font=\small] at (3,-0.8) {Taux de conformité $\mu$};
\node[font=\small, rotate=90] at (-0.9,2.5) {Taux d'incertitude $U$};
\end{tikzpicture}
\end{center}
\vspace{-0.4em}
\begin{center}{\small\itshape Figure~11.1. Le quadrant diagnostique et les
cinq domaines de LiZbiZ (d'après~\cite{minka2026nek}~; seuils du cours~;
échelle verticale $U \times 10$)}\end{center}
\end{minipage}

\medskip

\begin{tableaucours}{L{3.4cm}L{4.0cm}L{4.8cm}}{%
    \tableauHeader{Région} &
    \tableauHeader{Situation} &
    \tableauHeader{Action}
}
Conformité claire & On s'accorde, avec assurance, sur la conformité &
Validation courante \\
Non-conformité claire & On s'accorde, avec assurance, sur la
non-conformité & Correctifs \\
\textbf{Conformité opaque} & Des \OUI{} nombreux, mais une part importante
de \NEK{} & \textbf{Alerte prioritaire}~: rechercher la preuve avant tout
correctif \\
Non-conformité opaque & Des \NON{} nombreux et une part importante de
\NEK{} & Rechercher la preuve et engager des premiers correctifs \\
\end{tableaucours}
\vspace{-0.8em}
\begin{center}{\small\itshape Tableau~11.3. Les quatre régions et leurs
actions (d'après~\cite{minka2026nek})}\end{center}

La conformité opaque est la région la plus dangereuse~: une grille binaire
la confond avec la conformité claire, alors qu'elle peut cacher une
ignorance répandue. C'est exactement l'erreur du premier rapport de LiZbiZ.

\begin{boxinfo}
\textbf{Deux vocabulaires, une même idée.} L'article \emph{Vers
l'auditologie} distingue la \emph{non-conformité connue}, qui appelle des
correctifs, de la \emph{non-conformité inconnue}, qui appelle d'abord des
investigations~\cite{minka2026auditologie}. La première correspond à la
non-conformité claire~; la seconde aux deux régions opaques du quadrant.
\end{boxinfo}

% ============================================================================
\section{La Règle de Décision du Cours}
\label{sec:c11-regle}
% ============================================================================

\begin{boxprocedure}
\textbf{Conclure sur un domaine.} L'auditeur ne conclut qu'après avoir
consigné, dans le classeur O4, pour chaque domaine~:
\begin{enumerate}
    \item \textbf{$U$ d'abord.} Si $U \geq 0{,}20$, le domaine est opaque~:
          aucune conclusion de conformité n'est possible. L'auditeur
          recherche la preuve, ou conclut \emph{sur l'opacité elle-même}.
    \item \textbf{$\mu$ ensuite.} Si $U < 0{,}20$, le domaine est clair~:
          conforme si $\mu \geq 0{,}5$, non conforme sinon.
    \item \textbf{$H$ pour contrôler.} Une entropie élevée avec $U$ faible
          signale un désaccord entre \OUI{} et \NON{}~: le domaine n'est pas
          homogène (un site conforme, un autre non). On le décompose avant de
          conclure (Chapitre~\ref{chap:decomposition}). La décomposition de
          $H$ (section~\ref{sec:c11-decomp}) dit lequel des deux, ignorance
          ou désaccord, domine.
    \item \textbf{Les cotes de l'auditeur prévalent.} Les mesures des
          répondants orientent~; la conclusion repose sur les contrôles cotés
          par l'auditeur, preuve à l'appui (Chapitre~\ref{chap:preuve}).
\end{enumerate}
\end{boxprocedure}

\begin{boxlecture}{Les cinq domaines de LiZbiZ}
\begin{tableaucours}{L{4.0cm}cccL{3.2cm}}{%
    \tableauHeader{Domaine} &
    \tableauHeader{$\mu$} & \tableauHeader{$U$} & \tableauHeader{$H$} &
    \tableauHeader{Région}
}
A. Gouvernance du SI & 0,80 & 0,05 & 0,88 & Conformité claire \\
B. Sauvegardes & 0,67 & 0,08 & 1,19 & Conformité claire, $H$ élevé \\
C. Gestion des stocks & 0,08 & 0,08 & 0,82 & Non-conformité claire \\
D. Contrôle d'accès & 0,65 & 0,25 & 1,24 & \textbf{Conformité opaque} \\
E. Exploitation & 0,40 & 0,30 & 1,57 & Non-conformité opaque \\
\end{tableaucours}

\medskip
\textbf{Analyse} (confiance moyenne).
\begin{itemize}
    \item \emph{A}~: conclusion de conformité.
    \item \emph{B}~: conforme, mais l'entropie signale un désaccord~; la
          décomposition par site montre que l'agence de Garoua ne teste pas
          ses restaurations.
    \item \emph{C}~: non-conformité établie~: correctifs immédiats
          (Chapitre~\ref{chap:controle-interne}).
    \item \emph{D}~: alerte prioritaire~: une conformité apparente cache un
          quart de réponses inconnues~; preuves à rechercher avant toute
          conclusion.
    \item \emph{E}~: investigations et premiers correctifs.
\end{itemize}
\end{boxlecture}

\begin{boxalert}
\textbf{Violation typique.} Moyenner les cinq domaines pour donner une note
globale («~le SI de LiZbiZ est conforme à 52~\%~»).

\smallskip
\textbf{Règle~:} une moyenne efface les régions du quadrant. Le rapport
présente chaque domaine dans sa région~; s'il faut une synthèse chiffrée, ce
sont les comptes par région (deux domaines clairs conformes, un non
conforme clair, deux opaques).
\end{boxalert}

% ============================================================================
\section{L'Observation en Sept Éléments}
\label{sec:c11-observation}
% ============================================================================

\subsection{Sept éléments}

Chaque non-conformité significative donne lieu à une observation (outil~O8),
de structure invariante.

\begin{tableaucours}{cL{3.0cm}L{8.4cm}}{%
    \tableauHeader{N°} &
    \tableauHeader{Élément} &
    \tableauHeader{Contenu}
}
1 & Norme & Le critère, et le texte ou le référentiel d'où il vient \\
2 & Constatation & Le fait établi, chiffré, avec sa preuve et sa cote \\
3 & Risques & Ce qui peut arriver~: le risque de la fiche d'assertion et sa
conséquence \\
4 & Contradictoire & La réponse de l'audité, datée \\
5 & Analyse & La cause de l'écart et l'interprétation, avec un niveau de
confiance \\
6 & Conclusion & La position finale de l'auditeur \\
7 & Recommandations & Ce qu'il faut faire, en traitant la cause \\
\end{tableaucours}
\vspace{-0.8em}
\begin{center}{\small\itshape Tableau~11.4. Les sept éléments de
l'observation}\end{center}

\subsection{Sept éléments ou cinq~?}

Le cours MSN-2S3-MSI03 présente le constat d'audit en cinq éléments~:
constatation, critère, cause, conséquence, recommandation. Les sept
éléments les développent, sans les contredire.

\begin{tableaucours}{L{5.0cm}L{6.0cm}}{%
    \tableauHeader{Constat en cinq éléments} &
    \tableauHeader{Observation en sept éléments}
}
Critère & 1. Norme \\
Constatation & 2. Constatation \\
Conséquence & 3. Risques \\
(absent) & 4. Contradictoire \\
Cause & 5. Analyse \\
(absent) & 6. Conclusion \\
Recommandation & 7. Recommandations \\
\end{tableaucours}
\vspace{-0.8em}
\begin{center}{\small\itshape Tableau~11.5. Correspondance entre le constat
et l'observation}\end{center}

L'observation ajoute deux éléments que la méthode juge indispensables~: la
trace du contradictoire, qui rend l'observation opposable
(Chapitre~\ref{chap:mission}), et une conclusion explicite, qui sépare la
position de l'auditeur de son analyse.

\begin{boxobjet}{D.2.1. Observation}
\textbf{1. Norme~:} procédure interne de gestion des accès~: désactivation
du compte d'un agent parti sous 48~heures~; loi n°~2010/012, art.~26.

\textbf{2. Constatation~:} sur 31~départs entre janvier et avril~2026, 19
comptes ont été désactivés sous 48~heures, 7 plus tard, 5 sont encore
actifs (rapprochement exhaustif, papier de travail PT-14). Cote \NON{}.

\textbf{3. Risques~:} usage du compte d'un agent parti pour consulter ou
modifier des données de stock~; manquement à l'obligation de sécuriser le
système.

\textbf{4. Contradictoire~:} le responsable informatique indique, le
2~juillet, que les départs ne lui sont pas notifiés par les ressources
humaines.

\textbf{5. Analyse} (confiance élevée)~: la cause est l'absence de
notification des départs, non la négligence de l'informatique.

\textbf{6. Conclusion~:} le contrôle D.2.1 n'est pas efficace.

\textbf{7. Recommandations~:} notification automatique des départs par les
ressources humaines~; désactivation immédiate des cinq comptes~; revue
trimestrielle des comptes actifs.
\end{boxobjet}

\subsection{Graduer les observations}

Toutes les observations n'ont pas le même poids. Le cours retient trois
niveaux de gravité, à déclarer dans le plan d'audit.

\begin{tableaucours}{L{2.6cm}L{5.2cm}L{4.4cm}}{%
    \tableauHeader{Gravité} &
    \tableauHeader{Critère} &
    \tableauHeader{Exemple chez LiZbiZ}
}
Majeure & Risque résiduel intolérable, ou indice de fraude, ou manquement
légal & Ajustements de stock sans pièce~; TRANSLOG SERVICES \\
Significative & Contrôle inefficace sur un risque à traiter dans l'année &
D.2.1, cinq comptes d'agents partis actifs \\
Mineure & Défaut de forme, sans effet direct sur un risque & Compte rendu
de revue non daté \\
\end{tableaucours}
\vspace{-0.8em}
\begin{center}{\small\itshape Tableau~11.6. Les niveaux de gravité des
observations (convention du cours)}\end{center}

\subsection{Écrire une recommandation}

Une bonne recommandation traite la cause, pas le symptôme~; elle est
réalisable, vérifiable et attribuée.

\begin{tableaucours}{L{5.6cm}L{6.6cm}}{%
    \tableauHeader{Mal écrite} &
    \tableauHeader{Bien écrite}
}
Renforcer la sécurité des accès & Mettre en place une notification
automatique des départs des ressources humaines vers l'informatique, avant
le 31~août~2026 (responsable~: DRH) \\
Désactiver les comptes inutiles & Désactiver sous 48~heures les cinq
comptes d'agents partis identifiés (annexe~C), puis revoir chaque trimestre
les comptes actifs (responsable~: RSSI) \\
Licencier M.~FOUDA & (Hors du rôle de l'auditeur~: sanctionner relève de
la direction, après enquête.) \\
\end{tableaucours}
\vspace{-0.8em}
\begin{center}{\small\itshape Tableau~11.7. Écrire une recommandation}\end{center}

% ============================================================================
\section{Le Rapport}
\label{sec:c11-rapport}
% ============================================================================

\subsection{La structure}

\begin{boxprocedure}
\textbf{Le plan du rapport d'audit}
\begin{enumerate}
    \item \textbf{Synthèse} (une page)~: l'objectif, la conclusion par
          domaine, les observations majeures, ce qui reste inconnu.
    \item \textbf{Mandat et méthode}~: rappel de la lettre de mission,
          périmètre, référentiels, méthode (instrument \OUI{}, \NON{},
          \NEK{}, seuils du quadrant), limites et zones non auditées.
    \item \textbf{Lecture des domaines}~: pour chacun, $\mu$, $U$, $H$, la
          région du quadrant et la conclusion.
    \item \textbf{Observations}~: en sept éléments, classées par gravité.
    \item \textbf{Recommandations}~: hiérarchisées, avec responsable et
          échéance proposés.
    \item \textbf{Annexes}~: réponses écrites des audités, liste des
          documents reçus et manquants, méthode d'échantillonnage,
          manifeste des empreintes.
\end{enumerate}
\end{boxprocedure}

\subsection{La synthèse de LiZbiZ}

\begin{boxscenario}{Opération LIZBIZ. La synthèse du rapport (extrait)}
\textit{« La mission a évalué cinq domaines du système d'information de
LiZbiZ sur la période de janvier à avril~2026. La gouvernance (A) et les
sauvegardes (B) sont conformes, avec une réserve pour les restaurations de
l'agence de Garoua. La gestion des stocks (C) n'est pas conforme~: les
ajustements ne sont ni validés ni justifiés, pour 148~millions FCFA à
l'entrepôt de Douala. Le contrôle d'accès (D) et l'exploitation (E) ne
peuvent pas être jugés conformes~: un quart et un tiers des réponses
relèvent d'une connaissance insuffisante, faute de pièces~; la mission
recommande en priorité d'établir ces faits. Deux observations majeures
ont été portées à la connaissance du directeur général par note du
12~juin, conformément à la lettre de mission. La continuité d'activité n'a
pas été auditée. »}
\end{boxscenario}

\begin{boxalert}
\textbf{Violation typique.} Taire, dans la synthèse, les domaines opaques,
pour ne présenter que les non-conformités établies.

\smallskip
\textbf{Règle~:} ce que l'auditeur ne sait pas fait partie de sa conclusion.
Un domaine opaque est signalé comme tel, avec ce qu'il faudrait pour lever
l'incertitude.
\end{boxalert}

\subsection{Hiérarchiser les recommandations}

La région du quadrant donne l'ordre~: d'abord les observations majeures,
quelle que soit leur région~; puis les domaines en conformité opaque (où le
danger est caché)~; puis les non-conformités claires (où le correctif est
connu)~; enfin les améliorations des domaines conformes.

\begin{boxmetier}{Auditeur SI}
\textbf{Relire le rapport avant signature}
\begin{enumerate}
    \item Chaque chiffre renvoie-t-il à un papier de travail~?
    \item Chaque observation a-t-elle ses sept éléments, dont la réponse de
          l'audité~?
    \item La synthèse dit-elle ce qui n'est pas connu, et ce qui n'a pas été
          audité~?
    \item Le rapport contient-il une qualification pénale, un nom présenté
          comme coupable, une recommandation de sanction~? (Si oui, les
          retirer.)
    \item Les seuils du quadrant et de gravité sont-ils ceux du plan
          d'audit~?
\end{enumerate}
\end{boxmetier}

\begin{boxmetier}{Gestionnaire de risques}
\textbf{Intégrer les conclusions de l'audit}
\begin{enumerate}
    \item Reporter les non-conformités claires dans le registre des risques,
          comme risques dont les contrôles sont inefficaces.
    \item Traiter les domaines opaques comme un risque en soi~: on ne gère
          pas ce que l'on ne connaît pas.
    \item Suivre les recommandations dans le tableau de suivi
          (Chapitre~\ref{chap:suivi}).
\end{enumerate}
\end{boxmetier}

\begin{boxinfo}
\textbf{Ce que dit la recherche.} L'article A1~\cite{minka2026nek} propose
le quadrant et montre, sur des données d'audit, que des domaines jugés
conformes par une grille binaire se révèlent opaques avec \NEK{}. La
décomposition de l'entropie de ce chapitre en donne une lecture
supplémentaire~: elle sépare, dans la dispersion, ce qui relève de
l'ignorance et ce qui relève du désaccord.
\end{boxinfo}

\separateur

\begin{boxresume}
\textbf{Ce qu'il faut retenir de ce chapitre~:}
\begin{itemize}
    \item \textbf{Trois mesures}~: $\mu$ (part des \OUI{}), $U$ (part des
          \NEK{}), $H$ (dispersion)~; on agrège les réponses pour un domaine.
    \item \textbf{$H = H_b(U) + (1-U)\,H_b(\mu/(1-U))$}~: l'entropie se
          décompose en ignorance et désaccord~; deux entropies égales
          peuvent appeler deux démarches opposées.
    \item \textbf{Le quadrant} conformité × incertitude distingue
          conformité et non-conformité, claires ou opaques~; la conformité
          opaque est l'alerte prioritaire.
    \item \textbf{La règle de décision}~: $U$ d'abord, $\mu$ ensuite, $H$ pour
          contrôler~; jamais $H$ seul, jamais de moyenne globale.
    \item \textbf{L'observation en sept éléments} développe le constat en
          cinq~; elle est graduée~; la recommandation traite la cause.
    \item \textbf{Le rapport} dit aussi ce que l'auditeur ne sait pas et ce
          qu'il n'a pas audité.
\end{itemize}
\end{boxresume}

\begin{boxinfo}
\textbf{Points de contrôle de l'élève.} Je sais~: calculer $\mu$, $U$,
$H$~; démontrer la décomposition de $H$~; placer un domaine dans le
quadrant~; appliquer la règle de décision~; rédiger une observation
graduée~; écrire une recommandation qui traite la cause~; rédiger une
synthèse qui signale l'opacité.
\end{boxinfo}

\bigskip

\begin{boxexo}[title={\ding{45}\quad Exercice~11.1. Calculer et situer \niveaubase}]
Pour chacune des distributions suivantes (20 réponses), calculez $\mu$, $U$
et $H$, puis situez le domaine dans le quadrant.
\begin{enumerate}[(a)]
    \item 15 \OUI{}, 4 \NON{}, 1 \NEK{}.
    \item 6 \OUI{}, 4 \NON{}, 10 \NEK{}.
    \item 0 \OUI{}, 0 \NON{}, 20 \NEK{}.
    \item 3 \OUI{}, 16 \NON{}, 1 \NEK{}.
\end{enumerate}
\end{boxexo}

\begin{boxexo}[title={\ding{45}\quad Exercice~11.2. Décomposer \niveaumoyen}]
Décomposez l'entropie des distributions (a) et (b) de l'exercice~11.1, et
dites pour chacune si l'ignorance ou le désaccord domine. Quelle démarche
en déduisez-vous~?
\end{boxexo}

\begin{boxexo}[title={\ding{45}\quad Exercice~11.3. Démontrer un cas limite \niveaumoyen}]
À partir de la décomposition, montrez que $H = 0$ si et seulement si une
seule réponse est donnée par tous. Montrez ensuite que, pour $U$ fixé, $H$
est maximale lorsque $\mu = (1-U)/2$.
\end{boxexo}

\begin{boxexo}[title={\ding{45}\quad Exercice~11.4. Appliquer la règle \niveaumoyen}]
Appliquez la règle de décision aux domaines suivants de la BANQUE SANAGA~:
(a) changements 0,58/0,17~; (b) DIGICORE 0,50/0,13 avec $H = 1{,}41$~;
(c) continuité 0,20/0,10. Pour (b), que faites-vous de l'entropie~?
\end{boxexo}

\begin{boxexo}[title={\ding{45}\quad Exercice~11.5. Rédiger une observation \niveaumoyen}]
Rédigez l'observation en sept éléments relative aux ajustements de stock
de l'entrepôt de Douala (Chapitre~\ref{chap:fraude}), et indiquez sa
gravité.
\end{boxexo}

\begin{boxexo}[title={\ding{45}\quad Exercice~11.6. Corriger des recommandations \niveaumoyen}]
Réécrivez~: (a) «~Améliorer la gestion des sauvegardes~»~; (b) «~Former le
personnel à la sécurité~»~; (c) «~Mettre fin au tableur de Garoua~».
\end{boxexo}

\begin{boxexo}[title={\ding{45}\quad Exercice~11.7. Une note globale \niveaumoyen}]
Le directeur général de LiZbiZ demande « une note sur 20 » pour son conseil.
Que lui proposez-vous à la place, et pourquoi~? Rédigez le tableau que vous
remettez.
\end{boxexo}

\begin{boxexo}[title={\ding{45}\quad Exercice~11.8. Relire un rapport \niveaumoyen}]
La synthèse d'un rapport indique~: « Le contrôle d'accès est globalement
maîtrisé (65~\% de conformité). M.~KOUNOU porte la responsabilité des
défaillances. Nous recommandons son remplacement. » Listez les défauts et
réécrivez la synthèse.
\end{boxexo}

\begin{boxexo}[title={\ding{45}\quad Exercice~11.9. Deux domaines, une entropie \niveauexpert}]
Construisez deux distributions de 20 réponses qui ont la même entropie (à
0,02 près) mais des parts d'ignorance très différentes (moins de 20~\%
pour l'une, plus de 70~\% pour l'autre). Expliquez au comité d'audit, en
cinq lignes, pourquoi elles n'appellent pas la même action.
\end{boxexo}

\begin{boxexo}[title={\ding{45}\quad Exercice~11.10. Le rapport d'AUDIX CONSEIL \niveauexpert}]
À partir de la lecture des cinq domaines de LiZbiZ et des chapitres
précédents~:
\begin{enumerate}[(a)]
    \item Rédigez le plan détaillé du rapport, section par section.
    \item Hiérarchisez cinq recommandations en vous appuyant sur la
          gravité et le quadrant.
    \item Rédigez la section « Limites et zones non auditées ».
\end{enumerate}
\end{boxexo}

\bigskip

\begin{boxinfo}
\textbf{Pour aller plus loin.}
\begin{itemize}
    \item A1~\cite{minka2026nek}~: les mesures, le quadrant et leurs cas
          limites.
    \item C.~Shannon, \textit{A Mathematical Theory of
          Communication}~\cite{shannon1948}~: l'origine de l'entropie.
    \item T.~Cover et J.~Thomas, \textit{Elements of Information
          Theory}~\cite{cover2006elements}~: la règle de chaînage et la
          divergence de Kullback-Leibler, avec leurs démonstrations.
\end{itemize}
\end{boxinfo}

\begin{boxinfo}
\textbf{Transition.} Le rapport est remis. Mais l'audit ne s'arrête pas là~:
les recommandations doivent être suivies, et la mission suivante doit dire
si le système a changé. Le Chapitre~\ref{chap:suivi} présente le suivi, la
comparaison des cycles et la maturité.
\end{boxinfo}
